\chapter{Conclusions}
\label{ch:conclusions}

This thesis investigated how a commodity iOS device can support \textbf{interior design preview} through \textbf{parametric room capture}, \textbf{mesh-based editing}, optional \textbf{plenoptic visualisation}, \textbf{rule-based room analytics}, and \textbf{Blender export} for desktop refinement. The theoretical chapters argued that Apple's Object Capture instantiates a classical \textbf{photogrammetric} chain, whereas RoomPlan builds on \textbf{depth fusion} and learning to emit a structured \texttt{CapturedRoom}. The second theory chapter positioned \textbf{3D Gaussian Splatting} \cite{kerbl20233dgs} and feedforward \textbf{SHARP} as real-time alternatives to meshes for \emph{viewing}, not for CAD-style furniture replacement, and formalised a three-path comparison (Assets~A, B, and~C).

\section{Summary of Results}

The software artefact demonstrates a coherent pipeline: \textbf{RoomPlan} $\rightarrow$ serialised scan $\rightarrow$ \textbf{RealityKit} editor $\rightarrow$ USDZ replacements with gestures and decor $\rightarrow$ persistence across sessions. \textbf{Object Capture} extends asset authoring (Asset~A). A \textbf{Metal} splat renderer implements R-7 for imported \texttt{.ply} data (Asset~B). The \textbf{SHARP} integration (R-9, R-10) provides Asset~C via a local-network inference server without third-party cloud dependency for Apple capture paths. \textbf{Room analytics} (R-11) translate scan geometry into measurements, material estimates, and heuristic Smart Tips. \textbf{Blender export} (R-12) packages glTF structure, USDZ furniture, manifests, and a helper script so users can continue staging in desktop tools.

Requirements R-1--R-12 and associated non-functional goals N-1--N-8 are addressed at the design level in Chapter~\ref{ch:requirements} and at the code level in Chapter~\ref{ch:implementation}.

\section{Comparison of Reconstruction Modes}

Qualitative evaluation during development indicated that Asset~A (Object Capture meshes) delivers manipulable geometry aligned to RoomPlan bounding volumes and supports both in-app replacement and Blender export. Asset~B (offline-trained splats) achieves strong view-dependent realism when many calibrated views are available but does not support direct editing in the app. Asset~C (SHARP) minimises capture friction (single photograph, LAN inference after server warm-up) at the cost of single-view ambiguity and network dependency. None of the splat paths automatically register to \texttt{CapturedRoom} semantics; unified AR compositing would require future registration or fusion work.

\section{Limitations}

Accuracy is bounded by RoomPlan detection quality, photogrammetry assumptions in Object Capture, and the lack of closed-loop feedback between splats and parametric room semantics. Smart Tips are illustrative heuristics, not certified lighting design. Blender export reproduces the staged semantic scene, not photorealistic radiance fields; splats are excluded from the ZIP. 3DGS training remains external; real-time splat \emph{training} on device is research-grade \cite{Guo26}. Usability has not been quantified with SUS/HARUS in this thesis text; Kandil et al.\ \cite{kandil2021ux} provide a methodology for a future study.

\section{Future Work}

Promising extensions include: formal user evaluation; tighter integration between exported Object Capture image sets and external COLMAP/3DGS pipelines; coarse automatic registration of splats to room coordinates; optional inclusion of splat assets in DCC workflows; level-of-detail for very large \texttt{.ply} files; and semantic layout suggestions informed by work such as K\'an et al.\ \cite{kan2021ar}. Mithun et al.\ \cite{Geo2023} demonstrate that synthetic datasets with rich semantic labels can improve indoor scene understanding for AR applications, complementary to the real-scan approach taken here.

The project shows that \textbf{merging API-level room semantics with an editable mesh workflow and optional Blender hand-off} yields a practical middle ground between live AR furniture preview and professional BIM tools, suitable for iterative interior exploration on mobile hardware.
